\section{Open Questions}
Five substantive, testable questions remain after this replication. Each is
scoped for a self-contained follow-up pass on free endpoints and matched to
concrete next steps.

\subsection*{Q1. Epithelial re-fit (the ``Epi-'' isn't epithelial)}
Does the model retain descriptive power when re-fit against a genuine
epithelial system (BEAS-2B, HBEC-3KT, HMLE, MCF-10A, or a 3D airway/mammary
organoid $\gamma$-H2AX time-course) rather than confluent MRC-5
\emph{fibroblasts}? The paper's ``Epi-'' denotes \emph{epigenetic}, not
\emph{epithelial}; the entire parameter set is from Rothkamm \& L\"obrich 2003
lung fibroblasts. \textbf{Next steps:} digitise a public epithelial foci
curve; NLS-fit $\beta$, $B_T$, $\alpha$, $T$; report per-parameter shift and
residual RMSE. A $>2\times$ shift in $\beta$ would falsify cell-type
generality.

\subsection*{Q2. Explicit cell--cell coupling term}
Can the model be extended with a gap-junction or diffusible-bystander term
that its title advertises but its equations lack? \textbf{Next steps:}
introduce a coupling $\kappa$ into $\beta_{\text{eff}} = \beta(1 + \kappa
f_{\text{neighbors}})$; fit against paired confluent-vs-subconfluent or
2D-vs-3D coculture data; if $\kappa$ is statistically indistinguishable
from 0, the Epicellcom label is narrative-only; if non-zero, quantify how
much $\beta$ is inflated by the confluent-monolayer assumption.

\subsection*{Q3. EMT translation}
How does the model behave under epithelial-to-mesenchymal transition? EMT
modifies chromatin state, cell--cell junctions, and DDR pathway usage; none
are represented in Scott 2011. \textbf{Next steps:} find paired
epithelial/EMT $\gamma$-H2AX time-courses (e.g., MCF10A $\pm$ TGF-$\beta$);
fit the model to each state; test whether a single EMT-modulated
pathway-weight suffices to unify the fits.

\subsection*{Q4. Low-dose bystander threshold vs graded response}
Is the hard cut-off $T = 1.4$\,mGy defensible, or does reproducing
Rothkamm~\&~L\"obrich's own 5\,mGy hypersensitivity require a graded
low-dose response? \textbf{Next steps:} digitise the R\&L $\leq 20$\,mGy
foci curve; compare the threshold model against a smooth alternative
(e.g., BPM$(D) = \alpha D^2/(D+T)$ or a Hill transition); decide with
AICc/BIC.

\subsection*{Q5. Integration with organoid radiobiology}
Do 3D airway/intestinal organoid $\gamma$-H2AX resolution kinetics require
geometry-aware corrections (ROS-diffusion gradients, chromatin condensation)
that a 2D-monolayer-fit model cannot supply? Recent organoid DDR work
(Broutier 2017; Driehuis 2019; Bhaduri 2020) shows 2--3$\times$ slower
repair than matched 2D controls. \textbf{Next steps:} compile 3--5 paired
2D-vs-3D $\gamma$-H2AX datasets; fit Scott 2011 independently to each;
if a single geometry-scaling factor recovers 3D fits from 2D parameters,
propose it as a minimal Epicellcom-3D extension --- otherwise the model
should be documented as 2D-monolayer-limited.
